\chapter{1998:Maxim Kontsevich}\label{ch1998maximkontsevich:maxim-kontsevich-fields-medal-1998-a-comprehensive-report-on-his-work-and-impact}

\begin{center}\rule{0.5\linewidth}{0.5pt}\end{center}

\textbf{Prepared:} July 2026
\textbf{Subject:} Maxim Lvovich Kontsevich (born August 25, 1964, Khimki, USSR)
\textbf{Position:} Professor, IHES, Paris (since 1992)
\textbf{Doctoral Advisor:} Don Zagier (Universitat Bonn)
\textbf{Fields Medal Citation:} \emph{``For his contributions to algebraic
geometry, topology, and mathematical physics, in particular the
Witten conjecture.''}

\begin{center}\rule{0.5\linewidth}{0.5pt}\end{center}

\begin{prereqs}
Algebraic geometry, topology of moduli spaces, conformal field
theory, quantum field theory.
\end{prereqs}

\begin{objectives}
State Kontsevich's proof of the Witten conjecture; explain the KdV
hierarchy and intersection numbers on moduli spaces; describe the
Kontsevich matrix model.
\end{objectives}

% ============================================================
\section{Problem History}
% ============================================================

The \textbf{Witten conjecture} (1991) asserts that the generating
function of intersection numbers on the moduli space
$\overline{\mathcal{M}}_{g,n}$ of stable curves satisfies the
Korteweg--de Vries (KdV) hierarchy.  Let
\[
  F_g(t_0,t_1,\ldots)
  = \sum_{n=0}^{\infty}\frac{1}{n!}
    \int_{\overline{\mathcal{M}}_{g,n}}
    \psi_1^{d_1}\cdots\psi_n^{d_n}
\]
where $\psi_i$ is the first Chern class of the universal line
bundle and $d_1+\cdots+d_n=3g-3+n$.  Witten proposed that the
total potential $W=\sum_{g=0}^{\infty}\hbar^{g-1}F_g$ satisfies
the KdV hierarchy.

The \textbf{moduli space} $\overline{\mathcal{M}}_{g,n}$
parametrizes stable curves of genus $g$ with $n$ marked points.
It is a smooth Deligne--Mumford stack of dimension $3g-3+n$ (for
$2g-2+n>0$).  The \textbf{psi-class} $\psi_i$ on
$\overline{\mathcal{M}}_{g,n}$ is the first Chern class of the
cotangent line at the $i$-th marked point.  Intersection numbers
$\langle\tau_{d_1}\cdots\tau_{d_n}\rangle_g
 = \int_{\overline{\mathcal{M}}_{g,n}}\psi_1^{d_1}\cdots\psi_n^{d_n}$
are the fundamental invariants that Witten's conjecture relates to
integrable systems.

The \textbf{string equation}
$\langle\tau_0\tau_{d_1}\cdots\tau_{d_n}\rangle_g
 = \sum_{i=1}^{n}
   \langle\tau_{d_1}\cdots\tau_{d_i-1}\cdots\tau_{d_n}\rangle_g$
and the \textbf{dilaton equation}
$\langle\tau_1\tau_{d_1}\cdots\tau_{d_n}\rangle_g
 = (2g-2+n)\,\langle\tau_{d_1}\cdots\tau_{d_n}\rangle_g$
are fundamental recursion relations that follow from the geometry of
$\overline{\mathcal{M}}_{g,n}$.

The KdV hierarchy is a system of integrable PDEs:
\[
  \frac{\partial u}{\partial t_n}
  = \frac{\partial}{\partial x}R_n(u)
\]
where $R_n$ are polynomials in $u$ and its derivatives.  The first
equation $u_t = u_{xxx}+6uu_x$ is the KdV equation itself.

Before Kontsevich, Witten's conjecture was verified in low genera
by direct computation, and physicists provided heuristic arguments
using 2D quantum gravity.  But a rigorous mathematical proof was
lacking.

% ============================================================
\section{How Was It Solved}
% ============================================================

Kontsevich's proof proceeds in three steps.

\textbf{Step 1: Ribbon graphs.}
A \textbf{ribbon graph} is a graph embedded in a surface with
cyclic ordering of edges at each vertex.  The moduli space
$\mathcal{M}_{g,n}$ is homeomorphic to the space of ribbon graphs
of genus $g$ with $n$ legs, equipped with edge lengths.  Each edge
is assigned a length $\ell_e\in\mathbb{R}_{\geq 0}$, and the
resulting \textbf{metric ribbon graph} parameterizes Teichmuller
space.  A ribbon graph $\Gamma$ of genus $g$ with $n$ legs has
vertices of valence $\geq 3$, $n$ external edges (legs), and a
cyclic ordering of half-edges at each vertex such that the ribbon
surface has genus $g$.  The number of edges $e$ satisfies
$e=6g-6+3n$ for trivalent graphs.

The \textbf{Penner model} relates ribbon graphs to Weil--Petersson
volumes.  Kontsevich showed intersection numbers can be expressed
as:
\[
  \int_{\overline{\mathcal{M}}_{g,n}}
    \psi_1^{d_1}\cdots\psi_n^{d_n}
  = \sum_{\Gamma}\frac{1}{|\mathrm{Aut}(\Gamma)|}
    \prod_{\text{edges}}\ell_e
\]
where the sum is over ribbon graphs $\Gamma$.

\textbf{Step 2: Topological recursion relations.}
Intersection numbers satisfy a recurrence expressing integrals on
$\overline{\mathcal{M}}_{g,n+1}$ in terms of those on
$\overline{\mathcal{M}}_{g,n}$ and $\overline{\mathcal{M}}_{g-1,n+2}$.
This is the string equation and dilaton equation.

\textbf{Step 3: The matrix integral.}
The recurrence relations are equivalent to the KdV hierarchy via
the matrix integral:
\[
  Z = \int dM\,
       \exp\Bigl(-\tfrac{1}{2}\mathrm{tr}(\Lambda M^2)
                 +\hbar\,\mathrm{tr}(M^3)\Bigr)
\]
Expanding $Z$ in $\hbar$ gives the connected correlators, which are
exactly the intersection numbers.  The \textbf{Kontsevich matrix
model} is a Hermitian matrix integral where the action is
$N\,\mathrm{tr}(\Lambda M^2/2)-\mathrm{tr}(M^3)/3$.  The large-$N$
expansion produces Feynman diagrams that are precisely the ribbon
graphs.

% ============================================================
\section{Impact of This Discovery}
% ============================================================

Kontsevich's proof is one of the great achievements of modern
mathematical physics, establishing a rigorous connection between
string theory (matrix models) and algebraic geometry (intersection
numbers on moduli spaces).

The Kontsevich matrix model has become fundamental in random matrix
theory, combinatorics, and topological field theory.  It provides a
combinatorial interpretation of the KdV hierarchy.

The Witten conjecture influenced the study of \textbf{topological
recursion}, introduced by Eynard and Orantin as a generalization
computing invariants of algebraic curves, matrix models, and moduli
spaces.

Kontsevich's \textbf{homological mirror symmetry} conjecture
(1994) proposed that the derived category of coherent sheaves on one
Calabi--Yau manifold is equivalent to the Fukaya category of its
mirror.

The \textbf{Kontsevich deformation quantization} theorem (1999)
proved every Poisson manifold admits a deformation quantization,
settling a 1978 conjecture.

The proof also established the \textbf{Virasoro constraints} on the
generating function $F_g$, showing that the operators
$L_{-1},L_0,L_1,\dots$ from the Virasoro algebra annihilate
$e^F$.

% ============================================================
\section{Modern Views and Improvements}
% ============================================================

Topological recursion has been extended to a wide class of matrix
models and algebraic curves, providing a unified framework for
intersection numbers, Hurwitz numbers, and knot invariants.

Homological mirror symmetry has been proved for elliptic curves,
abelian varieties, and some Calabi--Yau threefolds, remaining a
major research area.

The Kontsevich matrix model has been generalized to multi-matrix
and non-commutative models, computing intersection numbers on more
general moduli spaces.

The \textbf{Eynard--Orantin topological recursion} has become a
universal tool for computing invariants across mathematics and
physics, from Hurwitz numbers to knot polynomials to
Gromov--Witten invariants of target spaces beyond curves.

Recent work by Andersen, Borot, and Orantin constructs the full
topological recursion from a matrix model using the \textbf{loop
equations}, providing a general framework that subsumes Kontsevich's
proof as a special case.

Modern work extends the Witten--Kontsevich hierarchy to higher
genera and to the Virasoro constraints.  Kontsevich himself
continues developing motivic integration, non-commutative
Grothendieck--Riemann--Roch, and quantum mechanics on moduli
spaces.

% ============================================================
\begin{center}\rule{0.5\linewidth}{0.5pt}\end{center}

\section*{Summary}
\addcontentsline{toc}{section}{Summary}

Kontsevich proved the Witten conjecture by showing the generating
function of intersection numbers on
$\overline{\mathcal{M}}_{g,n}$ satisfies the KdV hierarchy, using
the Kontsevich matrix model connecting ribbon graphs, Feynman
integrals, and moduli spaces.  He also proposed homological mirror
symmetry and proved the deformation quantization theorem.

% ============================================================
